\documentclass[10pt,a4paper]{amsart}
\usepackage[margin=19mm]{geometry}
\usepackage{amsmath,amssymb,mathtools}
\usepackage[scr=rsfs,cal=euler]{mathalfa}
\usepackage{xfrac}
\usepackage[unicode,hidelinks,bookmarks=false]{hyperref}
\newcommand{\ie}{i.e.\ }
\newcommand{\cf}{cf.\ }
\newcommand{\resp}{respectively\ }
\newcommand{\Subsection}{\subsection}
\providecommand{\nobreakdash}{\nobreak\hbox{-}}
\setlength{\emergencystretch}{2em}
\multlinegap0pt
\usepackage{amssymb}
\usepackage{bm}
\usepackage{mathtools}
\usepackage[shortlabels]{enumitem}
\newtheorem{proposition}{Proposition}
\newtheorem{theorem}{Theorem}
\title{One-Shot Optimization with Additional Inequality Constraints}
\author{Lea Fischer, Nicolas R. Gauger, Lisa Kusch}
\date{Source: arXiv:2606.02925v1 (CC BY 4.0)}
\begin{document}
\maketitle

\noindent\emph{Source and licence.} One-Shot Optimization with Additional Inequality Constraints. Version arXiv:2606.02925v1 (CC BY 4.0), distributed by arXiv under the Creative Commons Attribution 4.0 International licence, http://creativecommons.org/licenses/by/4.0/, as recorded in the arXiv licence metadata for this exact version and shown on the abstract page https://arxiv.org/abs/2606.02925v1. Full licence notice: \texttt{licenses/arxiv-2606.02925-CC-BY-4.0.txt}.\par\smallskip

\noindent\textit{Editorial scope.} Counted unit: the article's theorem thm:H1H2 (source0.tex lines 722--737). The article's complete original proof of it is delivered with it (source0.tex lines 739--781, one \texttt{proof} environment of 43 lines). No mathematics is written, repaired or re-derived: the statement and the proof are the article's own lines, in the article's own notation, and no retained line is substituted or re-flowed.  Retained uncounted as the source-local context the proof needs: lines 618--625; lines 666--680.  Nothing was omitted from any retained span, so the package carries no omission record.  No retained line contains a citation, so the wrapper carries no bibliography and no bibliography entry.  No retained line contains a figure environment, an image-inclusion command or drawing code, so no image is delivered and the package declares no asset dependency.  Editorial formatting only, in addition: an AMS \texttt{amsart} preamble that restates the article's own declarations and macros used by the retained lines, namely the environments \texttt{proposition}, \texttt{theorem} on separate counters, together with the standard packages those lines need.  The retained environments print in this document as Proposition 1, Theorem 1 in order of appearance, whereas the article prints the same items with the numbering of its own full document.  The complete native source of the article as distributed in the arXiv e-print for this exact version is bundled unchanged as \texttt{sources/201942/source0.tex}.
\begin{proposition}
\label{prophtwo}
    The active-set dependent contribution 
    \begin{align}
        \tilde H_2(D) = \Phi''_{2} = \begin{pmatrix} \alpha k_y^\top D k_y & \alpha k_y^\top D k_u & k_y^\top D \\[1ex] \alpha k_u^\top D k_y & \alpha k_u^\top D k_u & k_u^\top D \\[1ex] D k_y & D k_u & \frac1\alpha D \end{pmatrix}
    \end{align} 
    generated by the inequality term is positive semidefinite for every admissible generalized derivative selection $D(z)\in\partial\phi(z)$.
\end{proposition}

\begin{proposition}
\label{prop:H1}
Assume
\begin{align}
    \alpha\beta \Delta G_y^\top \Delta G_y \succ I+\beta N_{yy},
\end{align}
and
\begin{align}
    \alpha\beta k_yk_y^\top \succ I
\end{align}
Then
\[
\tilde H_1\succ0.
\]
\end{proposition}

\begin{theorem}
\label{thm:H1H2}

Assume that
\begin{align}
    \alpha\beta \Delta G_y^\top \Delta G_y \succ I+\beta N_{yy},
\end{align}
and
\begin{align}
    \alpha\beta k_yk_y^\top\succ I.
\end{align}

Then, for every admissible generalized derivative selection
$D(z)\in\partial\phi(z),$ the matrix $\tilde H_1+\tilde H_2(D)$ is positive definite.

\end{theorem}

\begin{proof}

By Proposition~\ref{prop:H1},
\[
\tilde H_1\succ0.
\]

Furthermore, Proposition~\ref{prophtwo} shows that
\[
\tilde H_2(D)\succeq0
\]
for every admissible generalized derivative selection
\[
D(z)\in\partial\phi(z).
\]

Since $\tilde H_1\succ0$, there exists a constant
$\gamma>0$ such that
\[
v^\top\tilde H_1v
\ge
\gamma\|v\|^2
\]
for all $v$.

Therefore,
\[
v^\top(\tilde H_1+\tilde H_2(D))v
=
v^\top\tilde H_1v
+
v^\top\tilde H_2(D)v
\ge
\gamma\|v\|^2.
\]

Hence
\[
\tilde H_1+\tilde H_2(D)\succ0.
\]


\end{proof}

\end{document}
